%
\subsection*{Tasas, soluciones y datos iniciales}
\label{CH01-B01}
En una EDO, la variable independiente suele ser el tiempo $t$; $y(t)$ es el
estado y los parámetros describen la ley adoptada. Si $y$ se mide en una
unidad $U$, $y'$ se mide en $U$/tiempo. Una solución satisface la ecuación
en todo un intervalo y, en un PVI, también el dato $y(t_0)=y_0$.
Compruébense ambos por sustitución. El orden es el de la mayor derivada.
La ecuación es lineal si la incógnita y sus derivadas aparecen en primer
grado, sin productos entre ellas, con coeficientes que dependen sólo de la
variable independiente.

Derivar una familia para eliminar constantes produce una ecuación que
debe comprobarse; puede admitir más soluciones que la familia de partida.
Recuérdense la regla de la cadena y los dominios de las funciones.
En $y'=ky$, $k$ es una tasa continua: en un tiempo $T$, el factor de
crecimiento es $e^{kT}$, no $1+kT$ salvo como aproximación cuando $|kT|$
es pequeño.

%
\subsection*{Separación, equilibrios y dominio}
\label{CH01-B02}
Si $y'=f(t)g(y)$, los ceros de $g$ dan soluciones constantes admisibles
cuando la ecuación está definida. En una región donde $g(y)\ne0$,
\[
 \int\frac{dy}{g(y)}=\int f(t)\,dt+C.
\]
El dato inicial fija la constante y ayuda a elegir la rama real. Antes de
dividir, anótese qué valores se excluyen y compruébense después en la
ecuación original. Son útiles las fracciones simples y
$\int dz/z=\log|z|+C$ en intervalos que no contienen cero.
El intervalo máximo debe contener el dato inicial y respetar tanto la
ecuación como la fórmula: no se atraviesa un polo ni se toma una raíz fuera
de su dominio real. Una expresión implícita también puede ser una solución.
Consulta: solución del ejercicio \exref{M2-E0103}.

%
\subsection*{Campo de pendientes y unicidad}
\label{CH01-B03}
Para $y'=F(t,y)$, el signo de $F$ indica ascenso o descenso y $F=0$
indica tangente horizontal. Una solución de equilibrio $y=c$ exige
$F(t,c)=0$ para todos los tiempos del intervalo; una tangente horizontal
aislada no basta.

Si $F$ es continua en un entorno de $(t_0,y_0)$, existe una solución local.
Si además $\partial F/\partial y$ es continua allí, la solución local es
única. Son criterios suficientes: cuando sus hipótesis fallan hay que
examinar la ecuación. Para demostrar no unicidad se necesitan soluciones
distintas con el mismo dato, no sólo el fallo de un criterio. En un empalme
de ramas deben coincidir valor y derivada, y la ecuación debe cumplirse
también en el punto de unión. Continuidad y límites laterales permiten
comprobarlo. Consulta: solución del ejercicio \exref{M2-E0191}.

%
\subsection*{Balance y final físico del modelo}
\label{CH01-B04}
El balance de una cantidad almacenada $Q$ tiene la forma
\[
 Q'=\text{entrada}-\text{salida}+\text{producción}-\text{consumo}.
\]
Todos sus términos deben tener unidades de cantidad por tiempo. Defínanse
estado, parámetros, signos e hipótesis antes de elegir una técnica.
Si el volumen depende de una altura, $V'=\frac{dV}{dh}h'$; esta regla de
la cadena permite cambiar de estado sin confundir altura y volumen.
El equilibrio se obtiene del balance, no de omitir un término incómodo.

Calibrar una constante con un dato no confirma la ley completa. Una
capacidad nominal tampoco fija el caudal real. Al vaciarse un recipiente o
alcanzarse un límite, el modelo puede terminar: continuar exige una regla
de borde o una nueva hipótesis física. Compruébense proporciones, geometría
y unidades. Consultas: soluciones de \exref{M2-E0025} y \exref{M2-E0032}.
